\clearpage
\SourceLeaf{189}
\thispagestyle{empty}
\pdfbookmark[0]{Anzeigen nicht eigner Schriften}{other-notices-title}
\begin{center}
\vspace*{1in}{\Large ANZEIGEN\par}
\vspace{.4in}{\large NICHT EIGNER\par}
\vspace{.4in}{\Large SCHRIFTEN.\par}
\end{center}

\clearpage
\SourceLeaf{191}
\renewcommand{\BookRunningTitle}{Anzeigen nicht eigner Schriften}
\pdfbookmark[1]{Recherches sur l'irréductibilité. Anzeige 1809 März 11}{notice-1809-03-11}
\begin{center}
\rule{.86\linewidth}{.3pt}\par
Göttingische gelehrte Anzeigen.\quad 1809 März 11.\par
\rule{.86\linewidth}{.3pt}
\end{center}
\begin{center}
\emph{Recherches sur l'irréductibilité Arithmétique et Géométrique des nombres et de leurs puissances.}\quad 1808.\quad (Ohne Druckort.\quad 25 S. in gr. Quart.)
\end{center}

Eine Schrift, deren Zweck dahin geht, die irrationalen Wurzelgrössen in Gestalt von rationalen Grössen darzustellen. Wir müssen uns begnügen, die Freunde der Mathematik auf dieses Werkchen aufmerksam gemacht zu haben, da die Grenzen dieser Blätter uns nicht verstatten, in die Darstellung und Prüfung des dem Verf. eigenthümlichen Gesichtspunkts und der von der gewöhnlichen ganz abgehenden Behandlung der Wurzelgrössen hier umständlicher einzugehen.

\Needspace{8\baselineskip}
\pdfbookmark[1]{Chernac, Cribrum Arithmeticum. Anzeige 1812 März 23}{notice-1812-03-23}
\begin{center}
\rule{.86\linewidth}{.3pt}\par
Göttingische Gelehrte Anzeigen.\quad 1812 März 23.\par
\rule{.86\linewidth}{.3pt}
\end{center}
\emph{Cribrum Arithmeticum, sive tabula continens numeros primos a compositis segregatos, occurrentes in serie numerorum ab unitate progredientium usque ad decies centena milia et ultra haec ad viginti millia} $(1020000)$. \emph{Numeris compositis, per} $2$, $3$, $5$ \emph{non dividuis, adscripti sunt divisores simplices, non minimi tantum, sed omnino}
\SourceLeaf{192}
\emph{omnes. Confecit} \textsc{Ladislaus Chernac}, \emph{Pannonius, A. L. M. Philos. et Medic. Doctor, in almo lyceo Daventriensi philosophiae professor. Daventriae} 1811. (Auf Kosten des Verfassers, gedruckt bei J. H. Lange.\quad XXII u. 1022 S. gr. Quart.)

Der vollständige Titel dieses wichtigen und sehr verdienstlichen Werks bezeichnet den Inhalt schon hinreichend: es ist eine durch eine eben so sorgfältige als mühsame Arbeit von mehreren Jahren berechnete Tafel für alle einfache Factoren aller durch $2$, $3$ und $5$ nicht theilbaren Zahlen von $1$ bis $1020000$, sauber und, soviel wir bei hin und wider angestellter Prüfung gefunden haben, sehr correct gedruckt. Wie schätzbar ein solches der Arithmetik gemachtes Geschenk sei, beurtheilt ein Jeder leicht, der viel mit grössern Zahlenrechnungen zu thun hat. Der Verf. verdient doppelten Dank, sowohl für seine höchst mühsame Arbeit selbst, wodurch er seinen Namen den unvergesslichen von \textsc{Rhaeticus}, \textsc{Pitiscus}, \textsc{Brigg}, \textsc{Vlacq}, \textsc{Wolfram}, \textsc{Taylor} u. A. zugesellt hat, als für den gewiss sehr erheblichen auf den Druck gemachten Aufwand, wofür sich sonst schwerlich ein Verleger gefunden haben möchte. Schon öfters sind dergleichen Tafeln, obwohl meistens in geringerer Ausdehnung, berechnet, aber entweder ganz im Manuscripte geblieben, oder im Abdruck nicht vollendet. \textsc{Lambert} munterte bekanntlich ehedem nach besten Kräften zur Fortsetzung der \textsc{Pell}schen, bis $100000$ gehenden und oft abgedruckten, Tafel auf, und einer von \textsc{Bernoulli} in \textsc{Lambert}'s Briefwechsel gegebenen Nachricht zufolge hatte \textsc{Oberreit} sie bis $500000$ fortgeführt, wovon die Abschrift in \textsc{Schulze}'s Hände gekommen war. \textsc{Anton Felkel} hatte sie, wie in der Monatl. Correspondenz 2. Bd. S. 223 berichtet wird, bis zu zwei Millionen in der Handschrift vollendet, und wollte sie späterhin bis $2460000$ geben; allein was davon in Wien auf öffentliche Kosten bereits gedruckt war, wurde, weil sich keine Käufer fanden, im Türkenkriege zu Patronen verbraucht! So ging eine verdienstliche vieljährige Arbeit für das Publicum verloren: um so mehr hielten wir es für Pflicht, die Erscheinung des gegenwärtigen Werks hier anzuzeigen. Die erste Million ist nun für Jedermanns Gebrauch da; und wer Gelegenheit und Eifer für diesen Gegenstand hat, möge daher seine Mühe auf das Weitere richten.

\SourceLeaf{193}
\Needspace{8\baselineskip}
\pdfbookmark[1]{Burckhardt, deuxième million. Anzeige 1814 November 3}{notice-1814-11-03}
\begin{center}
\rule{.86\linewidth}{.3pt}\par
Göttingische gelehrte Anzeigen.\quad 1814 November 3.\par
\rule{.86\linewidth}{.3pt}
\end{center}
\emph{Tables des diviseurs pour tous les nombres du deuxième million, ou plus exactement depuis} $1020000$ \emph{à} $2028000$, \emph{avec les nombres premiers qui s'y trouvent. Par} \textsc{J. Ch. Burckhardt}, \emph{membre de l'institut impérial, du bureau des longitudes de France, et de plusieurs autres sociétés savantes. Paris}, 1814. \emph{M\textsuperscript{me} V\textsuperscript{e} Courcier}. (VIII. u. 112 S. in Folio.)

Früher, als wir bei der Anzeige der die erste Million umfassenden Factorentafel von \textsc{Chernac} zu hoffen gewagt hätten, können wir schon die Vollendung und Erscheinung einer ähnlichen Tafel für die zweite Million berichten. Der verdiente Verfasser, dessen Name schon die grösste Sorgfalt und Genauigkeit verbürgt, hat sich durch diese mühsame Arbeit alle Freunde der Arithmetik sehr verpflichtet. \textsc{Chernac}'s Tafel für die erste Million gibt alle einfachen Factoren; die \textsc{Burckhardt}'sche für die zweite hingegen nur jedesmal den kleinsten Divisor. Die vollständige Zerlegung einer Zahl der zweiten Million erfordert also die Division mit dem kleinsten Divisor und das Aufsuchen des Quotienten in der \textsc{Chernac}'schen Tafel: allein diese kleine Mühe ist von gar keiner Erheblichkeit gegen den grossen Vortheil, die Tafel in einem so viel kleineren Raum zu besitzen, wobei die Aussicht bleibt, mit der Zeit die Tafel noch bis zu zehn Millionen ausgedehnt zu sehen. Die Zusammendrängung in den kleinen Band hat der Verfasser theils durch die Beschränkung auf den kleinsten Divisor, theils durch einen möglichst öconomischen Druck möglich gemacht. Wenn $a$ unbestimmt jede der achtzig Zahlen unter $300$ bedeutet, die durch $2$, $3$ und $5$ nicht theilbar sind, so ist überhaupt jede durch $2$, $3$ und $5$ nicht theilbare Zahl in der Form $300n+a$ begriffen. Alle achtzig Zahlen, für welche $n$ einerlei Werth hat, finden sich in Einer verticalen Columne, und solcher Columnen enthält jede Seite dreissig. Jede Seite umfasst also von neuntausend in der natürlichen Ordnung fortschreitenden Zahlen alle, welche durch $2$, $3$ oder $5$ nicht theilbar sind.

Die Methode, nach welcher Herr \textsc{Burckhardt} seine Tafel construirt hat, verdient hier noch eine besondere Erwähnung. Er liess ein Netz in Kupfer stechen, wo durch $81$ horizontale und $78$ verticale Linien ein in $80\times77$ d. i. $6160$ kleine Quadrate getheiltes Rechteck gebildet wurde, und davon die nöthige Anzahl von Abdrücken machen. An der Seite konnten sogleich die achtzig Werthe
\SourceLeaf{194}
von $a$ mit gestochen werden; die Werthe von $300n$ in fortlaufender Ordnung wurden mit der Feder über die $77$ verticalen Columnen geschrieben. So stellt jedes Blatt alle durch $2$, $3$ und $5$ nicht theilbaren Zahlen vor, welche unter je $23100$ in natürlicher Ordnung fortschreitenden Zahlen befindlich sind, und $44$ Blätter sind hinreichend, eine ganze Million zu umfassen. Man sieht leicht, dass die Zahlen, deren kleinster Theiler $7$ oder $11$ ist, auf jedem folgenden Blatte in derselben Ordnung wiederkehren, daher diese Divisoren sogleich auf die Kupferplatte gestochen werden konnten, und mithin auf jedem Blatte schon von selbst an den gehörigen Plätzen erschienen. Um nun die folgenden Divisoren z. B. $13$ einzutragen, nahm Herr B. von einem überzähligen Blatt der Breite nach blos $13$ Columnen, und indem er dasselbe als den Anfang seiner Tafel betrachtete, schnitt er alle die Quadrate, die den Divisor $13$ enthalten mussten, aus. Er brauchte also dieses Gitter nur auf die dreizehn ersten Columnen des ersten Blattes zu legen, dann auf die dreizehn folgenden u. s. w., um sogleich alle Plätze zu sehen, die, in so fern sie nicht schon $7$ oder $11$ enthielten, mit $13$ ausgefüllt werden mussten. Eben so wurde nachher mit dem Divisor $17$ u. s. w. verfahren. Bis zum Divisor $73$ reichten auf diese Weise die überzähligen Blätter hin; für die grössern Divisoren $79$, $83$ u. s. w. scheint Herr B. den Rahmen aus zwei oder mehreren Theilen zusammengesetzt zu haben. Bei den Divisoren hingegen, die über $500$ hinausgehen, zog Herr B. vor, die Vielfachen durch Addition zu suchen, wobei er für den andern Factor blos die Primzahlen zu nehmen brauchte. Wir finden dies ganze Verfahren höchst zweckmässig, und würden es allen denen zur Nachahmung empfehlen, die etwa Neigung haben sollten, die Tafel noch weiter fortzusetzen. Für die dritte und vierte Million hat inzwischen der Verfasser selbst schon einen grossen Theil der Rechnungen ausgeführt, daher wir gegründete Hoffnung haben, auch diese demnächst durch den Druck bekannt gemacht zu sehen.

\Needspace{8\baselineskip}
\pdfbookmark[1]{Burckhardt, troisième million. Anzeige 1816 November 7}{notice-1816-11-07}
\begin{center}
\rule{.86\linewidth}{.3pt}\par
Göttingische gelehrte Anzeigen.\quad 1816 November 7.\par
\rule{.86\linewidth}{.3pt}
\end{center}
\emph{Tables des diviseurs pour tous les nombres du troisième million, ou plus exactement, depuis} $2028000$ \emph{à} $3036000$, \emph{avec les nombres premiers qui s'y trouvent; par} \textsc{J. Chr. Burckhardt}, \emph{membre de l'académie royale des sciences, du bureau des }
\SourceLeaf{195}
\emph{longitudes de France et de plusieurs autres sociétés savantes. Paris} 1816. \emph{M\textsuperscript{me} V\textsuperscript{e} Courcier}. (112 Seiten in Folio.)

Da wir bereits bei der Anzeige der Tafel für die Factoren der zweiten Million die von dem verdienten Verf. angewandte Berechnungsmethode und die Einrichtung der Tafel selbst umständlich beschrieben haben, so können wir uns hier mit der blossen Anzeige von der Erscheinung der Tafel für die dritte Million begnügen. In Kurzem haben wir nun auch noch die Tafel für die \emph{erste} Million, auf dieselbe Art dargestellt von dem Verf. zu erwarten, so dass dann die ganze Tafel bis über drei Millionen nur einen mässigen Band ausmachen wird. Dem Verf. gebührt dafür der Dank aller Freunde der Arithmetik, die durch diese mühsame Arbeit ein Bedürfniss in einer Ausdehnung befriedigt sehen, die alles, was man noch vor wenigen Jahren zu hoffen wagen konnte, weit übersteigt.

\Needspace{8\baselineskip}
\pdfbookmark[1]{Burckhardt, premier million. Anzeige 1817 August 9}{notice-1817-08-09}
\begin{center}
\rule{.86\linewidth}{.3pt}\par
Göttingische gelehrte Anzeigen.\quad 1817 August 9.\par
\rule{.86\linewidth}{.3pt}
\end{center}
\emph{Tables des diviseurs, pour tous les nombres du premier million, ou plus exactement depuis} $1$ \emph{à} $1020000$, \emph{avec les nombres premiers qui s'y trouvent; par} \textsc{J. Chr. Burckhardt}, \emph{membre de l'académie des sciences dans l'institut royal, du bureau des longitudes de France, et de plusieurs autres sociétés savantes. Paris} 1817. \emph{M\textsuperscript{me} V\textsuperscript{e} Courcier}. (114 Seiten in Folio.)

Indem wir uns hier auf die Anzeigen der Tafeln für die zweite und dritte Million beziehen, kündigen wir jetzt blos das wirkliche\textperiodcentered\ Erscheinen dieser Factorentafeln für die erste Million an. Wir besitzen also nunmehr ein zusammenhängendes Ganzes für die drei ersten Millionen. Für die gegenwärtige erste Million bediente sich der Verfasser theils des \emph{Cribrum Arithmeticum} von \textsc{Chernac}, theils einer handschriftlichen Tafel von \textsc{Schenmark}, welche die Bibliothek des Königlichen Instituts besitzt. Letztere war indessen nicht ganz mit aller zu wünschenden Sorgfalt construirt, und die Entscheidung in Fällen, wo beide von einander abwichen, welche von beiden Recht habe, war oft ziemlich mühsam. In der \textsc{Chernac}'schen Tafel zeigte sich nur eine sehr geringe Anzahl von Fehlern, welche Herr \textsc{Burckhardt} hier mitgetheilt hat.

\SourceLeaf{196}
Auch für die vierte, fünfte und sechste Million hat der Verf. die Materialien bereits grösstentheils vorräthig, und er erbietet sich, diese Fortsetzung zu liefern, wenn der Verleger durch einen hinreichenden Absatz der drei ersten Millionen aufgemuntert wird. Es wäre in der That sehr zu beklagen, wenn die Früchte einer so mühsamen und nützlichen Arbeit der Welt entzogen werden sollten.
